\section*{Введение}
\addcontentsline{toc}{section}{Введение}

Напомним (\rf{1.44}), что поле $\Om$ \emph{алгебраически замкнуто}, если каждый непостоянный многочлен из $\Om[X]$ имеет
корень в $\Om$ (тогда он раскладывается в $\Om[X]$ на линейные множители), и что \emph{алгебраическое замыкание}
поля $F$~--- это алгебраически замкнутое поле $\Om\supset F$, алгебраическое над $F$. В этой главе мы докажем, что
алгебраическое замыкание есть у каждого поля и что любые два алгебраических замыкания поля $F$ изоморфны над $F$.
В двух случаях существование устанавливается просто.

\begin{itemize}[leftmargin=*]
\item \emph{Подполя $\CC$.} Если $F\subset\CC$, то множество $\{z\in\CC\mid z\ \text{алгебраичен над}\ F\}$~--- поле
(\rf{1.46}) и алгебраическое замыкание $F$ (\rf{1.47}; здесь используется алгебраическая замкнутость $\CC$, \rf{5.6}). Оно задано без
всякого произвола; что это означает, обсуждается в \rf{6.8a}.
\item \emph{Счётные поля.} Пусть $F$ не более чем счётно, т.\,е.\ конечно или счётно (например, $F=\F_p$, $\Q$ или $\Q(T)$). Многочлен задаётся конечным
набором коэффициентов, поэтому унитарные неприводимые многочлены из $F[X]$ можно перенумеровать: $f_1,f_2,\dots$
Положим $E_0=F$ и $E_i$~--- поле разложения $f_i$ над $E_{i-1}$. Объединение цепочки
$E_0\subset E_1\subset E_2\subset\cdots$ (понимаемое как в \rf{4.25}: поле разложения строится абстрактно, и $E_{i-1}$
отождествляется со своим образом в $E_i$)~--- поле $\Om$, алгебраическое над $F$: каждый его элемент лежит в
некотором $E_i$, а $[E_i:F]<\infty$. Каждый многочлен из $F[X]$~--- произведение константы и нескольких $f_i$, и потому
раскладывается в $\Om[X]$ на линейные множители; по \rf{1.45} поле $\Om$ алгебраически замкнуто. Похожим образом
(объединением цепочки полей разложения) в \rf{4.24} было построено $\Fpb$.
\end{itemize}

В общем случае трудность~--- в теории множеств. Хочется взять <<объединение полей разложения всех многочленов из
$F[X]$>>, но это нужно сделать способом, который допускают аксиомы теории множеств: совокупность всех алгебраических
расширений поля $F$ (если $F$ не алгебраически замкнуто) не является множеством (как и совокупность всех множеств), и строить замыкание <<внутри неё>>
нельзя. Ниже мы напомним лемму Цорна и нужные факты о мощностях, а затем приведём три способа обойти эту трудность.

Несчётные поля, не лежащие в $\CC$, встречаются естественно. Например, поле рядов Лорана $\F_p((T))$ имеет
характеристику $p$ и несчётно: уже степенных рядов $\sum_{i\ge0}a_iT^i$ столько же, сколько последовательностей
$(a_i)$ элементов $\F_p$, а их несчётно много:
$|\F_p^{\mathbb N}|\ge|\{0,1\}^{\mathbb N}|=|\mathcal P(\mathbb N)|>\aleph_0$ (\rf{6.4d}).
